\section*{Resumen}
\addcontentsline{toc}{section}{Resumen}

El proceso de pre-sustentación de trabajos de titulación en la Universidad Técnica Estatal de Quevedo
(UTEQ) se gestionaba con trámites manuales y hojas de cálculo dispersas, sin trazabilidad verificable
entre solicitud, jurado, cronograma, evaluación y acta. Este trabajo presenta el diseño, implementación y
evaluación empírica de un sistema web (Spring Boot 3.2.1, Angular 21, PostgreSQL 15, Redis 7) que
automatiza ese flujo, siguiendo Design Science Research de Peffers et al. El sistema implementa 15
requisitos funcionales (12 con prueba automatizada real) con matriz de trazabilidad bidireccional. La
evaluación empírica combina estadística descriptiva e inferencial real: 5 corridas de carga k6 (p95 de
6.17--8.53\,ms, 50 usuarios concurrentes), test de Mann-Whitney comparando caché fría/caliente de Redis
(latencia 13.9$\times$ menor, $p<0.001$), auditoría de seguridad OWASP Top~10, y seis corridas de
Lighthouse. El desarrollo estuvo marcado por hallazgos reales corregidos explícitamente: datos de
usabilidad y trazabilidad previamente fabricados fueron detectados y reemplazados por mediciones reales.
Los resultados muestran un sistema funcional con brechas documentadas honestamente: los 8 requisitos Must
tienen prueba dedicada (100\,\%), pero 3 de 15 aún no (80\,\% general), y la cobertura de código llega a
63.17\,\% de líneas (395 pruebas), declarado como trabajo futuro. Una revisión manual del modelo de
autorización encontró y corrigió 6 fallos reales de control de acceso (IDOR y mass-assignment) invisibles
al análisis automatizado. Este informe documenta resultados positivos y limitaciones reales, priorizando
la integridad de la evidencia sobre la completitud aparente.

\textbf{Palabras clave:} ingeniería de requisitos; Design Science Research; evaluación empírica de
software; seguridad de aplicaciones web; usabilidad; integridad de la investigación.

\textbf{Categorías CCS (ACM Computing Classification System 2012):}
Software and its engineering~Requirements analysis;
Software and its engineering~Empirical software validation;
Security and privacy~Web application security;
Information systems~Web applications.

\newpage

\section*{Abstract}
\addcontentsline{toc}{section}{Abstract}

The pre-defense process for undergraduate thesis projects at Universidad Técnica Estatal de Quevedo
(UTEQ) was managed through manual paperwork and scattered spreadsheets, with no verifiable traceability
between requests, juries, schedules, evaluations, and minutes. This work presents the design,
implementation, and empirical evaluation of a web system (Spring Boot 3.2.1, Angular 21, PostgreSQL 15,
Redis 7) that automates this workflow, following Peffers et al.'s Design Science Research methodology.
The system implements 15 functional requirements (12 with real automated test coverage) with a
bidirectional traceability matrix. Empirical evaluation combines real descriptive and inferential
statistics: 5 k6 load-testing runs (p95 of 6.17--8.53\,ms, 50 concurrent users), a Mann-Whitney test
comparing cold and warm Redis cache states (13.9$\times$ lower latency, $p<0.001$), an OWASP Top~10
security audit, and six Lighthouse runs. Development was marked by real findings explicitly corrected
rather than hidden: previously fabricated usability and traceability data were detected and replaced
with real measurements. Results show a functional system with honestly documented gaps: all 8
Must-priority requirements have a dedicated test (100\,\%), but 3 of 15 still lack one (80\,\% overall),
and code coverage reaches 63.17\,\% of lines (395 tests), declared as future work. A manual review of the
authorization model found and fixed 6 real access-control flaws (IDOR and mass-assignment) invisible to
automated analysis. This report documents positive results and real limitations, prioritizing evidence
integrity over apparent completeness.

\textbf{Keywords:} requirements engineering; Design Science Research; empirical software evaluation; web
application security; usability; research integrity.
\newpage
